\section{Job Sequencing Problem}
\label{sec:job-sequencing}

\problemstatement{We are given $n$ jobs, where the $i$-th job has a deadline
$\textit{deadline}[i]$ and a profit $\textit{profit}[i]$. Every job takes
exactly one unit of time to complete, and only one job can be processed at
any given unit of time (time slots are $1, 2, 3, \dots$). A job earns its
profit only if it is completed at or before its deadline. We must choose
which jobs to schedule, and at which time slots, so as to \textbf{maximize
total profit}; we are also asked to report how many jobs end up being
scheduled.}

\begin{patternbox}
\emph{``Each job has a deadline and a profit; schedule jobs (each taking one
unit of time) to maximize total profit''} is the signature of the
\emph{Job Sequencing Problem}. The two keywords to notice together are
\textbf{``profit, maximize''} (suggesting we should prioritize high-value
jobs first) and \textbf{``deadline''} (suggesting that, once we have decided
\emph{which} jobs to take, we still need a rule for placing them into time
slots without violating their deadlines). This combination -- rank by value,
then place greedily as late as possible -- is a different flavor of greedy
from interval scheduling, and is often paired conceptually with the Disjoint
Set Union (Union-Find) data structure for an even faster slot-finding step.
\end{patternbox}

\subsection*{Understanding the problem carefully}

Two separate questions are tangled together here: \emph{which} jobs should
we choose to do, and \emph{when} should we do each chosen job. Profit is
earned per job, independent of which legal slot it occupies, so among
several candidate sets of jobs we might choose, more total profit is always
better, with no benefit to a particular schedule beyond it being legal at
all. This suggests considering jobs in decreasing order of profit, and
greedily trying to fit each one in, since a job with higher profit is always
worth taking over a lower-profit job if there is room for only one of them.

Once we decide to take a particular job with deadline $d$, the only real
question is \emph{which} slot among $1, 2, \dots, d$ to place it in. Here a
second, equally important greedy idea applies: place the job in the
\textbf{latest} available slot not exceeding its deadline, rather than the
earliest. The reasoning is that an early slot is more valuable to keep
\emph{open}, because it can still accommodate jobs with very tight (small)
deadlines that arrive later in our profit-sorted processing order, whereas a
late slot close to $d$ is useless to any job whose deadline is smaller than
$d$ anyway.

\begin{ideabox}[Greedy Choice]
Sort all jobs in decreasing order of profit. For each job, in this order,
scan backward from its deadline toward slot $1$ and place it in the first
free slot found. If no free slot exists at or before its deadline, skip this
job permanently -- it can never be scheduled later, since slots only get
more occupied as we proceed, not less.
\end{ideabox}

\subsection*{Why processing by profit, and placing as late as possible, is correct}

Consider why we never want to revisit a profit-order decision. Once a
higher-profit job has successfully claimed a slot, no later (lower-profit)
job can ever displace it without decreasing total profit, so committing
immediately, in profit order, is safe -- this is the same "highest value
first" exchange argument used throughout greedy scheduling: if some optimal
solution skipped a higher-profit job $J$ to make room for a lower-profit job
$J'$ that conflicts with it, swapping $J$ in for $J'$ (which is always
possible, since $J$ has at least as early a deadline requirement satisfied
by any slot $J'$ could use, or an even more flexible range) cannot decrease,
and may increase, total profit.

Placing each accepted job as late as possible within its own deadline window
is the secondary greedy choice that makes the slot-assignment feasible
without any backtracking. By reserving early slots for as long as possible,
we never block a future job (processed later, because it has lower profit)
that might desperately need one of those early slots due to having a very
small deadline -- a job with deadline $1$ has nowhere else to go, so an
early slot occupied unnecessarily by an earlier, more flexible job would be
a wasted opportunity that greedy avoids by always preferring the latest
legal slot first.

\subsection*{Algorithm}

\begin{enumerate}
  \item Sort the jobs in decreasing order of profit.
  \item Let $\textit{maxDeadline}$ be the largest deadline among all jobs.
        Create a boolean array $\textit{slot}[1..\textit{maxDeadline}]$,
        all initially free.
  \item Initialize $\textit{countJobs} \gets 0$ and
        $\textit{totalProfit} \gets 0$.
  \item For each job, in decreasing profit order, with deadline $d$:
  \begin{enumerate}
    \item Scan $t$ from $\min(d, \textit{maxDeadline})$ down to $1$:
    \begin{enumerate}
      \item If $\textit{slot}[t]$ is free: mark it occupied, increment
            $\textit{countJobs}$, add this job's profit to
            $\textit{totalProfit}$, and stop scanning (break out of this
            inner loop) -- the job has been placed.
    \end{enumerate}
    \item If no free slot was found in the scan, this job is permanently
          skipped.
  \end{enumerate}
  \item Return $\textit{countJobs}$ and $\textit{totalProfit}$.
\end{enumerate}

\begin{algorithm}[H]
\caption{Job Sequencing}
\begin{algorithmic}[1]
\Procedure{JobSequencing}{\textit{jobs}}
  \State sort \textit{jobs} in decreasing order of profit
  \State $\textit{maxDeadline} \gets \max_i \textit{deadline}[i]$
  \State $\textit{slot}[1..\textit{maxDeadline}] \gets$ all free
  \State $\textit{countJobs} \gets 0$, $\textit{totalProfit} \gets 0$
  \For{each job $(d, p)$ in sorted order}
    \For{$t \gets \min(d, \textit{maxDeadline})$ \textbf{down to} $1$}
      \If{$\textit{slot}[t]$ is free}
        \State mark $\textit{slot}[t]$ occupied
        \State $\textit{countJobs} \gets \textit{countJobs} + 1$
        \State $\textit{totalProfit} \gets \textit{totalProfit} + p$
        \State \textbf{break}
      \EndIf
    \EndFor
  \EndFor
  \State \Return $(\textit{countJobs}, \textit{totalProfit})$
\EndProcedure
\end{algorithmic}
\end{algorithm}

\subsection*{Dry run}

\dryrun{Take jobs with (deadline, profit) pairs:
$(4,20), (1,10), (1,40), (1,30)$.}

Sorted by decreasing profit:
$(1,40), (1,30), (4,20), (1,10)$.

Slots $1, 2, 3, 4$ all start free.

\begin{itemize}
  \item Job $(1,40)$: scan from $t=1$ down to $1$. Slot $1$ is free; place
        it there. $\textit{countJobs}=1$, $\textit{totalProfit}=40$.
        Slots: $[1{=}\text{occupied}, 2,3,4{=}\text{free}]$.
  \item Job $(1,30)$: scan from $t=1$ down to $1$. Slot $1$ is occupied.
        No free slot found within its deadline. Skip this job.
  \item Job $(4,20)$: scan from $t=4$ down to $1$. Slot $4$ is free; place
        it there. $\textit{countJobs}=2$, $\textit{totalProfit}=60$.
        Slots: $[1{=}\text{occ}, 2,3{=}\text{free}, 4{=}\text{occ}]$.
  \item Job $(1,10)$: scan from $t=1$ down to $1$. Slot $1$ is occupied.
        Skip.
\end{itemize}

The final result is $\textit{countJobs} = 2$ jobs scheduled, with
$\textit{totalProfit} = 60$.

\subsection*{C++ implementation}

\begin{lstlisting}
#include <bits/stdc++.h>
using namespace std;

struct Job {
    int id;
    int deadline;
    int profit;
};

pair<int,int> jobSequencing(vector<Job>& jobs) {
    sort(jobs.begin(), jobs.end(), [](const Job& a, const Job& b) {
        return a.profit > b.profit;
    });

    int maxDeadline = 0;
    for (auto& job : jobs) {
        maxDeadline = max(maxDeadline, job.deadline);
    }

    vector<bool> slot(maxDeadline + 1, false); // 1-indexed
    int countJobs = 0, totalProfit = 0;

    for (auto& job : jobs) {
        for (int t = min(job.deadline, maxDeadline); t >= 1; t--) {
            if (!slot[t]) {
                slot[t] = true;
                countJobs++;
                totalProfit += job.profit;
                break;
            }
        }
    }
    return {countJobs, totalProfit};
}

int main() {
    vector<Job> jobs = {
        {1, 4, 20},
        {2, 1, 10},
        {3, 1, 40},
        {4, 1, 30}
    };

    auto result = jobSequencing(jobs);
    cout << "Jobs done: " << result.first << endl;
    cout << "Total profit: " << result.second << endl;
    return 0;
}
\end{lstlisting}

\begin{complexitybox}
\textbf{Time Complexity.} Sorting the $n$ jobs by profit costs
$O(n \log n)$. For each job, scanning backward through slots costs up to
$O(\textit{maxDeadline})$ in the worst case, giving a total of
$O(n \cdot \textit{maxDeadline})$ for the slot-placement phase. Hence the
overall time complexity is
\[
  O(n \log n + n \cdot \textit{maxDeadline}).
\]
This can be improved to $O(n \log n + n \log n) = O(n \log n)$ overall by
using a Disjoint Set Union structure (covered in Chapter~6) to jump directly
to the latest free slot at or before a given deadline, rather than scanning
slot by slot.
\textbf{Space Complexity.} We use a boolean array of size
$O(\textit{maxDeadline})$ to track slot occupancy, so the space complexity
is $O(\textit{maxDeadline})$, in addition to $O(n)$ for sorting the jobs.
\end{complexitybox}

\begin{takeawaybox}
When a problem combines \emph{``maximize total value''} with
\emph{``each item also has a deadline/capacity constraint on where it can
be placed,''} a two-layer greedy often applies: sort by value to decide
\emph{which} items to take, and within that order, place each accepted item
as late (or otherwise as flexibly) as possible so as not to block items that
have less placement freedom and are still waiting to be considered.
\end{takeawaybox}
